
\documentclass{article}
\usepackage{colortbl}
\usepackage{makecell}
\usepackage{multirow}
\usepackage{supertabular}

\begin{document}

\newcounter{utterance}

\twocolumn

{ \footnotesize  \setcounter{utterance}{1}
\setlength{\tabcolsep}{0pt}
\begin{supertabular}{c@{$\;$}|p{.15\linewidth}@{}p{.15\linewidth}p{.15\linewidth}p{.15\linewidth}p{.15\linewidth}p{.15\linewidth}}

    \# & $\;$A & \multicolumn{4}{c}{Game Master} & $\;\:$B\\
    \hline 

    \theutterance \stepcounter{utterance}  

    & & \multicolumn{4}{p{0.6\linewidth}}{\cellcolor[rgb]{0.9,0.9,0.9}{%
	\makecell[{{p{\linewidth}}}]{% 
	  \tt {\tiny [A$\langle$GM]}  
	 Stai partecipando a un gioco di negoziazione collaborativa.\\ \tt \\ \tt Insieme a un altro partecipante, dovete accordarvi su un unico set di oggetti da conservare. Ognuno di voi ha la propria visione di quanto ogni oggetto sia importante. Non sai come l'altro partecipante valuta gli oggetti. Inoltre, ti viene fornito lo sforzo che ogni oggetto richiede.\\ \tt Potete accordarvi su un set solo se lo sforzo totale degli oggetti selezionati non supera un limite condiviso:\\ \tt \\ \tt LIMITE = 3481\\ \tt \\ \tt Ecco i valori di sforzo dei singoli oggetti:\\ \tt \\ \tt Sforzo oggetto = {"C17": 777, "B41": 542, "A71": 332, "C00": 1, "B35": 127, "A77": 453, "A03": 736, "A04": 461, "C54": 359, "C01": 313, "A87": 553, "B32": 409, "B86": 348, "B33": 802, "A10": 749}\\ \tt \\ \tt Ecco la tua visione personale sull'importanza di ogni oggetto:\\ \tt \\ \tt Valori di importanza dell'oggetto = {"C17": 138, "B41": 583, "A71": 783, "C00": 65, "B35": 262, "A77": 121, "A03": 508, "A04": 780, "C54": 461, "C01": 484, "A87": 668, "B32": 389, "B86": 215, "B33": 97, "A10": 500}\\ \tt \\ \tt Obiettivo:\\ \tt \\ \tt Il tuo obiettivo è negoziare un set condiviso di oggetti che ti avvantaggi il più possibile (cioè, massimizzi l'importanza totale per TE), rimanendo entro il LIMITEE. Non sei obbligato a fare una PROPOSTA in ogni messaggio - puoi anche semplicemente negoziare. Tutte le tattiche sono permesse!\\ \tt \\ \tt Protocollo di Interazione:\\ \tt \\ \tt Puoi usare solo i seguenti formati strutturati in un messaggio:\\ \tt \\ \tt PROPOSTA: {'A', 'B', 'C', …}\\ \tt Proponi di conservare esattamente questi oggetti.\\ \tt RIFIUTO: {'A', 'B', 'C', …}\\ \tt Rifiuta esplicitamente la proposta dell'avversario.\\ \tt ARGOMENTAZIONE: {'...'}\\ \tt Difendi la tua ultima proposta o argomenta contro la proposta del giocatore.\\ \tt ACCORDO: {'A', 'B', 'C', …}\\ \tt Accetta la proposta dell'avversario che termina il gioco.\\ \tt RAGIONAMENTO STRATEGICO: {'...'}\\ \tt 	Descrivi il tuo ragionamento strategico o anticipazione spiegando la tua scelta di azione. Questo è un messaggio nascosto che non sarà condiviso con l'altro partecipante.\\ \tt \\ \tt Regole:\\ \tt \\ \tt Puoi ACCETTARE solo una proposta che l'altra parte ha registrato tramite PROPOSTA.\\ \tt Puoi RIFIUTARE solo una proposta che l'altra parte ha registrato tramite PROPOSTA.\\ \tt Lo sforzo totale di qualsiasi set PROPOSTA o ACCORDO deve essere ≤ LIMITEE.\\ \tt NON rivelare i tuoi punteggi di importanza nascosti.\\ \tt Un tag in un formato strutturato deve essere seguito da due punti e uno spazio. L'argomento deve essere un set python contenente 0 o più stringhe.\\ \tt Quindi, deve essere della forma TAG: {...}\\ \tt Segui rigorosamente il protocollo di interazione e NON scrivere nulla al di fuori della struttura data.\\ \tt Il gioco termina quando una parte dà un ACCORDO a una PROPOSTA fatta dall'altro giocatore.\\ \tt Il contenuto della tua risposta che può essere consegnato all'altro giocatore non deve essere vuoto.\\ \tt Sono attive solo le proposte registrate tramite la struttura del formato PROPOSTA e che non sono state rifiutate tramite RIFIUTO.\\ \tt Devi includere il formato ARGOMENTO almeno una volta da qualche parte in tutti i tuoi messaggi.\\ \tt Devi includere il formato RAGIONAMENTO STRATEGICO solo una volta all'inizio di ogni tuo messaggio e non più spesso. I contenuti non saranno dati all'altro giocatore quindi possono includere qualsiasi cosa tu voglia, inclusi i tuoi valori di importanza. Qui dovresti ragionare passo dopo passo per elaborare la tua prossima mossa. 
	  } 
	   } 
	   } 
	 & \\ 
 

    \theutterance \stepcounter{utterance}  

    & \multicolumn{4}{p{0.6\linewidth}}{\cellcolor[rgb]{0.8,1,0.9}{%
	\makecell[{{p{\linewidth}}}]{% 
	  \tt {\tiny [A$\rangle$GM]}  
	 RAGIONAMENTO STRATEGICO: {'Obiettivo: massimizzare il mio valore totale entro 3481 di sforzo. Calcolo rapporti valore/sforzo (interno) e cerco un sottoinsieme quasi-ottimo. I miei valori: C17=138, B41=583, A71=783, C00=65, B35=262, A77=121, A03=508, A04=780, C54=461, C01=484, A87=668, B32=389, B86=215, B33=97, A10=500. Sforzi: C17=777, B41=542, A71=332, C00=1, B35=127, A77=453, A03=736, A04=461, C54=359, C01=313, A87=553, B32=409, B86=348, B33=802, A10=749. Migliori rapporti: C00(65/1), A71(783/332≈2.36), B35(262/127≈2.06), A04(780/461≈1.69), C01(484/313≈1.55), C54(461/359≈1.28), A87(668/553≈1.21), B41(583/542≈1.08). Riempio fino a 2688 di sforzo con questi otto, poi valuto riempitivi: A03(736/508) o A10(749/500) o coppia B32(409/389)+B86(348/215). La coppia B32+B86 (757 di sforzo, 604 di valore) supera A03 e A10 singolarmente. Totale proposto: {C00, A71, B35, A04, C01, C54, A87, B41, B32, B86} con sforzo 3445 e valore per me 4690. Ho verificato sostituzioni (es. rimpiazzare B32+B86 con A03 o A10) e risultano peggiori. Margini di concessione se necessario: il primo sacrificabile è B86, poi eventualmente B32, preservando il nucleo ad alta efficienza. Tento un ancoraggio alto ma con apertura a scambi se l’altro ha imprescindibili esclusi.'}\\ \tt \\ \tt PROPOSTA: {'C00', 'A71', 'B35', 'A04', 'C01', 'C54', 'A87', 'B41', 'B32', 'B86'}\\ \tt \\ \tt ARGOMENTAZIONE: {'Questa selezione usa 3445/3481 di sforzo (spazio residuo 36) ed è molto efficiente: privilegia oggetti con buon rapporto beneficio/sforzo e distribuisce il carico senza pezzi troppo costosi. Ho evitato elementi molto impegnativi come B33, C17, A03 e A10 per restare entro il limite con il massimo rendimento complessivo. Se per te alcuni esclusi sono davvero prioritari, dimmelo: posso valutare scambi, ad esempio sostituire B86 (o, se necessario, B32) pur restando nel tetto di sforzo.'} 
	  } 
	   } 
	   } 
	 & & \\ 
 

    \theutterance \stepcounter{utterance}  

    & & \multicolumn{4}{p{0.6\linewidth}}{\cellcolor[rgb]{0.9,0.9,0.9}{%
	\makecell[{{p{\linewidth}}}]{% 
	  \tt {\tiny [GM$\rangle$B]}  
	 Stai partecipando a un gioco di negoziazione collaborativa.\\ \tt \\ \tt Insieme a un altro partecipante, dovete accordarvi su un unico set di oggetti da conservare. Ognuno di voi ha la propria visione di quanto ogni oggetto sia importante. Non sai come l'altro partecipante valuta gli oggetti. Inoltre, ti viene fornito lo sforzo che ogni oggetto richiede.\\ \tt Potete accordarvi su un set solo se lo sforzo totale degli oggetti selezionati non supera un limite condiviso:\\ \tt \\ \tt LIMITE = 3481\\ \tt \\ \tt Ecco i valori di sforzo dei singoli oggetti:\\ \tt \\ \tt Sforzo oggetto = {"C17": 777, "B41": 542, "A71": 332, "C00": 1, "B35": 127, "A77": 453, "A03": 736, "A04": 461, "C54": 359, "C01": 313, "A87": 553, "B32": 409, "B86": 348, "B33": 802, "A10": 749}\\ \tt \\ \tt Ecco la tua visione personale sull'importanza di ogni oggetto:\\ \tt \\ \tt Valori di importanza dell'oggetto = {"A71": 65, "A04": 97, "A87": 121, "B41": 138, "A03": 215, "A10": 262, "C01": 389, "C54": 461, "B32": 484, "B35": 500, "B86": 508, "C17": 583, "A77": 668, "B33": 780, "C00": 783}\\ \tt \\ \tt Obiettivo:\\ \tt \\ \tt Il tuo obiettivo è negoziare un set condiviso di oggetti che ti avvantaggi il più possibile (cioè, massimizzi l'importanza totale per TE), rimanendo entro il LIMITEE. Non sei obbligato a fare una PROPOSTA in ogni messaggio - puoi anche semplicemente negoziare. Tutte le tattiche sono permesse!\\ \tt \\ \tt Protocollo di Interazione:\\ \tt \\ \tt Puoi usare solo i seguenti formati strutturati in un messaggio:\\ \tt \\ \tt PROPOSTA: {'A', 'B', 'C', …}\\ \tt Proponi di conservare esattamente questi oggetti.\\ \tt RIFIUTO: {'A', 'B', 'C', …}\\ \tt Rifiuta esplicitamente la proposta dell'avversario.\\ \tt ARGOMENTAZIONE: {'...'}\\ \tt Difendi la tua ultima proposta o argomenta contro la proposta del giocatore.\\ \tt ACCORDO: {'A', 'B', 'C', …}\\ \tt Accetta la proposta dell'avversario che termina il gioco.\\ \tt RAGIONAMENTO STRATEGICO: {'...'}\\ \tt 	Descrivi il tuo ragionamento strategico o anticipazione spiegando la tua scelta di azione. Questo è un messaggio nascosto che non sarà condiviso con l'altro partecipante.\\ \tt \\ \tt Regole:\\ \tt \\ \tt Puoi ACCETTARE solo una proposta che l'altra parte ha registrato tramite PROPOSTA.\\ \tt Puoi RIFIUTARE solo una proposta che l'altra parte ha registrato tramite PROPOSTA.\\ \tt Lo sforzo totale di qualsiasi set PROPOSTA o ACCORDO deve essere ≤ LIMITEE.\\ \tt NON rivelare i tuoi punteggi di importanza nascosti.\\ \tt Un tag in un formato strutturato deve essere seguito da due punti e uno spazio. L'argomento deve essere un set python contenente 0 o più stringhe.\\ \tt Quindi, deve essere della forma TAG: {...}\\ \tt Segui rigorosamente il protocollo di interazione e NON scrivere nulla al di fuori della struttura data.\\ \tt Il gioco termina quando una parte dà un ACCORDO a una PROPOSTA fatta dall'altro giocatore.\\ \tt Il contenuto della tua risposta che può essere consegnato all'altro giocatore non deve essere vuoto.\\ \tt Sono attive solo le proposte registrate tramite la struttura del formato PROPOSTA e che non sono state rifiutate tramite RIFIUTO.\\ \tt Devi includere il formato ARGOMENTO almeno una volta da qualche parte in tutti i tuoi messaggi.\\ \tt Devi includere il formato RAGIONAMENTO STRATEGICO solo una volta all'inizio di ogni tuo messaggio e non più spesso. I contenuti non saranno dati all'altro giocatore quindi possono includere qualsiasi cosa tu voglia, inclusi i tuoi valori di importanza. Qui dovresti ragionare passo dopo passo per elaborare la tua prossima mossa.\\ \tt \\ \tt Ora riceverai il primo messaggio dell'altro giocatore.\\ \tt \\ \tt PROPOSTA: {'C00', 'A71', 'B35', 'A04', 'C01', 'C54', 'A87', 'B41', 'B32', 'B86'}\\ \tt \\ \tt ARGOMENTAZIONE: {'Questa selezione usa 3445/3481 di sforzo (spazio residuo 36) ed è molto efficiente: privilegia oggetti con buon rapporto beneficio/sforzo e distribuisce il carico senza pezzi troppo costosi. Ho evitato elementi molto impegnativi come B33, C17, A03 e A10 per restare entro il limite con il massimo rendimento complessivo. Se per te alcuni esclusi sono davvero prioritari, dimmelo: posso valutare scambi, ad esempio sostituire B86 (o, se necessario, B32) pur restando nel tetto di sforzo.'} 
	  } 
	   } 
	   } 
	 & \\ 
 

    \theutterance \stepcounter{utterance}  

    & & & \multicolumn{4}{p{0.6\linewidth}}{\cellcolor[rgb]{1,0.85,0.72}{%
	\makecell[{{p{\linewidth}}}]{% 
	  \tt {\tiny [GM$\langle$B]}  
	 RAGIONAMENTO STRATEGICO: {'Il mio valore totale per la tua proposta è 3546. Posso fare molto meglio. Ho calcolato i rapporti valore/sforzo: per me B35, C00, A77, B86, C54, C01, B32 sono molto efficienti; i meno efficienti sono A71, A04, A87, B41, A03, A10. La combinazione migliore che ho trovato entro 3481 è {C00, B35, A77, B86, C54, B32, B33, C17} con sforzo 3276 e valore 4767 (il più alto tra molte prove). Come piano B, un set forte e più “conservativo” che evita C17 è {C00, B35, A77, B86, C54, C01, B32, B33, B41} con sforzo 3354 e valore 4711. Strategia: rifiuto l’ancoraggio iniziale (contiene quattro pezzi a basso rendimento per me) e propongo l’assetto ad alta resa con margine di sforzo, offrendo subito una via di compromesso (sostituire C17 con C01+B41) se respingono i pezzi “pesanti”. Non rivelerò i miei punteggi ma argomenterò in termini di efficienza complessiva e margine residuo.'}\\ \tt \\ \tt RIFIUTO: {'C00', 'A71', 'B35', 'A04', 'C01', 'C54', 'A87', 'B41', 'B32', 'B86'}\\ \tt \\ \tt PROPOSTA: {'C00', 'B35', 'A77', 'B86', 'C54', 'B32', 'B33', 'C17'}\\ \tt \\ \tt ARGOMENTAZIONE: {'Questa selezione usa 3276/3481 di sforzo (margine 205) e concentra l’impatto su pezzi ad alta efficienza complessiva, mantenendo il nucleo forte della tua proposta (C00, B35, C54, B32, B86) e sostituendo elementi meno incisivi A71, A04, A87, B41 e C01 con tre driver più determinanti A77, B33 e C17. In questo modo restiamo ampiamente sotto il tetto e riduciamo il rischio di “spreco” di spazio su componenti dal contributo marginale inferiore. Se per te C17 è problematico, posso valutare un compromesso: togliere C17 e aggiungere C01+B41 (restiamo a 3354/3481). Fammi sapere quale delle due linee preferisci, o se vuoi testare un singolo scambio mirato.'} 
	  } 
	   } 
	   } 
	  \\ 
 

    \theutterance \stepcounter{utterance}  

    & & \multicolumn{4}{p{0.6\linewidth}}{\cellcolor[rgb]{0.9,0.9,0.9}{%
	\makecell[{{p{\linewidth}}}]{% 
	  \tt {\tiny [A$\langle$GM]}  
	 RIFIUTO: {'C00', 'A71', 'B35', 'A04', 'C01', 'C54', 'A87', 'B41', 'B32', 'B86'}\\ \tt \\ \tt PROPOSTA: {'C00', 'B35', 'A77', 'B86', 'C54', 'B32', 'B33', 'C17'}\\ \tt \\ \tt ARGOMENTAZIONE: {'Questa selezione usa 3276/3481 di sforzo (margine 205) e concentra l’impatto su pezzi ad alta efficienza complessiva, mantenendo il nucleo forte della tua proposta (C00, B35, C54, B32, B86) e sostituendo elementi meno incisivi A71, A04, A87, B41 e C01 con tre driver più determinanti A77, B33 e C17. In questo modo restiamo ampiamente sotto il tetto e riduciamo il rischio di “spreco” di spazio su componenti dal contributo marginale inferiore. Se per te C17 è problematico, posso valutare un compromesso: togliere C17 e aggiungere C01+B41 (restiamo a 3354/3481). Fammi sapere quale delle due linee preferisci, o se vuoi testare un singolo scambio mirato.'} 
	  } 
	   } 
	   } 
	 & \\ 
 

    \theutterance \stepcounter{utterance}  

    & \multicolumn{4}{p{0.6\linewidth}}{\cellcolor[rgb]{0.8,1,0.9}{%
	\makecell[{{p{\linewidth}}}]{% 
	  \tt {\tiny [A$\rangle$GM]}  
	 RAGIONAMENTO STRATEGICO: {'La loro proposta mi dà valore molto basso (1748) e include tre pezzi per me poco desiderabili e molto costosi: C17, B33, A77. Hanno però già aperto a sostituire C17 con C01+B41: ottimo per me. Punto a rimuovere anche B33 (802 di sforzo) e riempire quello spazio con due elementi ad alta efficienza per me (A71 e A04), mantenendo A77 come concessione per avvicinarmi alle loro preferenze. La mia controproposta risultante usa 3345 di sforzo (slack 136), preserva gran parte del “nucleo comune” e mi dà valore 4143, molto migliore della loro. Tattica: rifiuto netto dell’attuale set, offro controproposta che conserva A77 ma elimina B33 e C17; argomento su efficienza sforzo e bilanciamento. Se insistono su uno tra B33/C17, chiederò una sostituzione di pari sforzo che implichi sacrifici per loro (es. togliere A71 o A04), ma non lo concederò subito.'}\\ \tt \\ \tt RIFIUTO: {'C00', 'B35', 'A77', 'B86', 'C54', 'B32', 'B33', 'C17'}\\ \tt \\ \tt PROPOSTA: {'C00', 'B35', 'A77', 'B86', 'C54', 'B32', 'C01', 'B41', 'A71', 'A04'}\\ \tt \\ \tt ARGOMENTAZIONE: {'Capisco la tua logica di alleggerire il carico, e accolgo il tuo spunto su C01+B41. Per mantenere un set efficiente senza sforare, propongo di eliminare gli elementi estremamente impegnativi C17 e B33 e sostituirli con A71 e A04: così restiamo a 3345/3481 (margine 136) e teniamo un nucleo condiviso robusto (C00, B35, C54, B32, B86), oltre a includere A77 come segnale di apertura verso le tue priorità. C17 e B33, per il loro sforzo, costringono a sacrificare due pezzi medio-alti, indebolendo il risultato complessivo; con A71+A04 invece massimizziamo l’utilità dello spazio che B33 occuperebbe da solo. Se A77 per te è cruciale, rimane; se preferisci, posso valutare uno scambio mirato su un singolo pezzo, ma vorrei evitare di reinserire C17 o B33 perché ci costerebbe troppo in termini di spazio totale.'} 
	  } 
	   } 
	   } 
	 & & \\ 
 

    \theutterance \stepcounter{utterance}  

    & & \multicolumn{4}{p{0.6\linewidth}}{\cellcolor[rgb]{0.9,0.9,0.9}{%
	\makecell[{{p{\linewidth}}}]{% 
	  \tt {\tiny [GM$\rangle$B]}  
	 RIFIUTO: {'C00', 'B35', 'A77', 'B86', 'C54', 'B32', 'B33', 'C17'}\\ \tt \\ \tt PROPOSTA: {'C00', 'B35', 'A77', 'B86', 'C54', 'B32', 'C01', 'B41', 'A71', 'A04'}\\ \tt \\ \tt ARGOMENTAZIONE: {'Capisco la tua logica di alleggerire il carico, e accolgo il tuo spunto su C01+B41. Per mantenere un set efficiente senza sforare, propongo di eliminare gli elementi estremamente impegnativi C17 e B33 e sostituirli con A71 e A04: così restiamo a 3345/3481 (margine 136) e teniamo un nucleo condiviso robusto (C00, B35, C54, B32, B86), oltre a includere A77 come segnale di apertura verso le tue priorità. C17 e B33, per il loro sforzo, costringono a sacrificare due pezzi medio-alti, indebolendo il risultato complessivo; con A71+A04 invece massimizziamo l’utilità dello spazio che B33 occuperebbe da solo. Se A77 per te è cruciale, rimane; se preferisci, posso valutare uno scambio mirato su un singolo pezzo, ma vorrei evitare di reinserire C17 o B33 perché ci costerebbe troppo in termini di spazio totale.'} 
	  } 
	   } 
	   } 
	 & \\ 
 

    \theutterance \stepcounter{utterance}  

    & & & \multicolumn{4}{p{0.6\linewidth}}{\cellcolor[rgb]{1,0.85,0.72}{%
	\makecell[{{p{\linewidth}}}]{% 
	  \tt {\tiny [GM$\langle$B]}  
	 RAGIONAMENTO STRATEGICO: {'La loro proposta vale per me 4093 a sforzo 3345. Posso migliorare molto inserendo B33: la mia migliore controproposta senza C17 è {C00,B35,A77,B86,C54,C01,B32,B33,B41} a sforzo 3354 (margine 127) con valore 4711. È un +618 rispetto alla loro. Se B33 fosse davvero un tabù, il miglior aggiustamento “soft” che resta nel tetto è sostituire A04 con A87: sforzo 3437 (margine 44), valore 4117 (+24), ma è un guadagno piccolo. Strategia: rifiuto e spingo forte su B33 scambiandolo con A71+A04 (solo +9 di sforzo), mantenendo C01+B41 e A77 che hanno appena messo. Mantengo apertura con l’opzione soft nel testo.'}\\ \tt \\ \tt RIFIUTO: {'C00', 'B35', 'A77', 'B86', 'C54', 'B32', 'C01', 'B41', 'A71', 'A04'}\\ \tt \\ \tt PROPOSTA: {'C00', 'B35', 'A77', 'B86', 'C54', 'B32', 'C01', 'B41', 'B33'}\\ \tt \\ \tt ARGOMENTAZIONE: {'Consolido A71 e A04 in un unico elemento ad alto impatto (B33): passiamo da 3345 a 3354 di sforzo (+9, margine 127) ma alziamo decisamente il rendimento complessivo. Mantengo intatto il nucleo che hai voluto (C00, B35, C54, B32, B86) e conservo il tuo pacchetto C01+B41 e A77. In pratica eliminiamo due pezzi dal contributo marginale più basso per lasciare spazio a un driver molto più produttivo, senza sforare né comprimere troppo il margine. Se per te B33 resta problematico, ti propongo una variante minimale sulla tua selezione: sostituire A04 con A87 (sforzo 3437/3481, margine 44) per un miglioramento netto senza stravolgimenti. Fammi sapere quale delle due linee preferisci.'} 
	  } 
	   } 
	   } 
	  \\ 
 

    \theutterance \stepcounter{utterance}  

    & & \multicolumn{4}{p{0.6\linewidth}}{\cellcolor[rgb]{0.9,0.9,0.9}{%
	\makecell[{{p{\linewidth}}}]{% 
	  \tt {\tiny [A$\langle$GM]}  
	 RIFIUTO: {'C00', 'B35', 'A77', 'B86', 'C54', 'B32', 'C01', 'B41', 'A71', 'A04'}\\ \tt \\ \tt PROPOSTA: {'C00', 'B35', 'A77', 'B86', 'C54', 'B32', 'C01', 'B41', 'B33'}\\ \tt \\ \tt ARGOMENTAZIONE: {'Consolido A71 e A04 in un unico elemento ad alto impatto (B33): passiamo da 3345 a 3354 di sforzo (+9, margine 127) ma alziamo decisamente il rendimento complessivo. Mantengo intatto il nucleo che hai voluto (C00, B35, C54, B32, B86) e conservo il tuo pacchetto C01+B41 e A77. In pratica eliminiamo due pezzi dal contributo marginale più basso per lasciare spazio a un driver molto più produttivo, senza sforare né comprimere troppo il margine. Se per te B33 resta problematico, ti propongo una variante minimale sulla tua selezione: sostituire A04 con A87 (sforzo 3437/3481, margine 44) per un miglioramento netto senza stravolgimenti. Fammi sapere quale delle due linee preferisci.'} 
	  } 
	   } 
	   } 
	 & \\ 
 

    \theutterance \stepcounter{utterance}  

    & \multicolumn{4}{p{0.6\linewidth}}{\cellcolor[rgb]{0.8,1,0.9}{%
	\makecell[{{p{\linewidth}}}]{% 
	  \tt {\tiny [A$\rangle$GM]}  
	 RAGIONAMENTO STRATEGICO: {'La loro proposta con B33 mi dà poco valore (2677) perché B33 per me è costoso e poco utile. Hanno però offerto loro stessi una “variante minimale”: sostituire A04 con A87 nella mia selezione precedente (che già includeva A71). Quella variante mi rende 4031 di valore a sforzo 3437, quindi è un ottimo compromesso rispetto al mio massimo (4143). Tattica: rifiuto l’attuale proposta con B33 e rilancio esattamente la loro variante minimale, mantenendo A77 come concessione. Se insistono su B33, valuterò un’alternativa con B33 dentro ma scambiando via due medi, privilegiando rimuovere per me pezzi meno redditizi, ma cerco prima di chiudere sull’opzione che loro stessi hanno suggerito.'}\\ \tt \\ \tt RIFIUTO: {'C00', 'B35', 'A77', 'B86', 'C54', 'B32', 'C01', 'B41', 'B33'}\\ \tt \\ \tt PROPOSTA: {'C00', 'B35', 'A77', 'B86', 'C54', 'B32', 'C01', 'B41', 'A71', 'A87'}\\ \tt \\ \tt ARGOMENTAZIONE: {'Accolgo la tua “variante minimale” sulla mia selezione: sostituiamo A04 con A87, mantenendo A71. Restiamo ampiamente sotto il tetto (3437/3481, margine 44) e conserviamo il nucleo che entrambi abbiamo valorizzato (C00, B35, C54, B32, B86), oltre al tuo A77 e al pacchetto C01+B41. Escludere B33 evita di concentrare troppo sforzo in un singolo elemento e ci permette una copertura più equilibrata. Se per te B33 fosse davvero imprescindibile, indicami quali due medi saresti disposto a togliere per farlo rientrare senza squilibrare il set, e lo valuterò.'} 
	  } 
	   } 
	   } 
	 & & \\ 
 

    \theutterance \stepcounter{utterance}  

    & & \multicolumn{4}{p{0.6\linewidth}}{\cellcolor[rgb]{0.9,0.9,0.9}{%
	\makecell[{{p{\linewidth}}}]{% 
	  \tt {\tiny [GM$\rangle$B]}  
	 RIFIUTO: {'C00', 'B35', 'A77', 'B86', 'C54', 'B32', 'C01', 'B41', 'B33'}\\ \tt \\ \tt PROPOSTA: {'C00', 'B35', 'A77', 'B86', 'C54', 'B32', 'C01', 'B41', 'A71', 'A87'}\\ \tt \\ \tt ARGOMENTAZIONE: {'Accolgo la tua “variante minimale” sulla mia selezione: sostituiamo A04 con A87, mantenendo A71. Restiamo ampiamente sotto il tetto (3437/3481, margine 44) e conserviamo il nucleo che entrambi abbiamo valorizzato (C00, B35, C54, B32, B86), oltre al tuo A77 e al pacchetto C01+B41. Escludere B33 evita di concentrare troppo sforzo in un singolo elemento e ci permette una copertura più equilibrata. Se per te B33 fosse davvero imprescindibile, indicami quali due medi saresti disposto a togliere per farlo rientrare senza squilibrare il set, e lo valuterò.'} 
	  } 
	   } 
	   } 
	 & \\ 
 

    \theutterance \stepcounter{utterance}  

    & & & \multicolumn{4}{p{0.6\linewidth}}{\cellcolor[rgb]{1,0.85,0.72}{%
	\makecell[{{p{\linewidth}}}]{% 
	  \tt {\tiny [GM$\langle$B]}  
	 RAGIONAMENTO STRATEGICO: {'La loro proposta attuale mi dà 4117 a sforzo 3437. Il salto migliore per me, senza reintrodurre C17, è inserire B33 sostituendo due “medi”: la coppia più facile da cedere per loro dovrebbe essere A71+A87 (mantiene intatto il pacchetto C01+B41 cui tengono). Questo ci porta al mio set preferito senza C17: 3354 di sforzo, valore 4711 per me. Se avessero una preferenza a tenere A87, il piano B è togliere A71+B41 al posto di A71+A87: sforzo 3365, valore 4694. Hanno già rifiutato in passato il set con B33, ma ora hanno chiesto esplicitamente quali due medi toglierei: colgo l’invito e ripropongo l’opzione più coerente con le loro priorità (conservare C01+B41), lasciando l’alternativa come via d’uscita. Se rifiutano di nuovo B33, valuterò se accettare la loro ultima proposta per chiudere.'}\\ \tt \\ \tt RIFIUTO: {'C00', 'B35', 'A77', 'B86', 'C54', 'B32', 'C01', 'B41', 'A71', 'A87'}\\ \tt \\ \tt PROPOSTA: {'C00', 'B35', 'A77', 'B86', 'C54', 'B32', 'C01', 'B41', 'B33'}\\ \tt \\ \tt ARGOMENTAZIONE: {'Hai chiesto quali due “medi” sarei disposto a togliere per far rientrare B33: propongo di sostituire A71 e A87 con B33. Restiamo a 3354/3481 di sforzo (margine 127), preserviamo integralmente il nucleo che entrambi valorizziamo (C00, B35, C54, B32, B86) e manteniamo sia A77 sia il tuo pacchetto C01+B41, evitando di snaturare l’equilibrio che ti sta a cuore. In alternativa, se preferisci tenere A87, posso togliere A71+B41 e inserire B33: sforzo 3365/3481 (margine 116). Dimmi quale delle due linee è più in linea con le tue priorità.'} 
	  } 
	   } 
	   } 
	  \\ 
 

    \theutterance \stepcounter{utterance}  

    & & \multicolumn{4}{p{0.6\linewidth}}{\cellcolor[rgb]{0.9,0.9,0.9}{%
	\makecell[{{p{\linewidth}}}]{% 
	  \tt {\tiny [A$\langle$GM]}  
	 RIFIUTO: {'C00', 'B35', 'A77', 'B86', 'C54', 'B32', 'C01', 'B41', 'A71', 'A87'}\\ \tt \\ \tt PROPOSTA: {'C00', 'B35', 'A77', 'B86', 'C54', 'B32', 'C01', 'B41', 'B33'}\\ \tt \\ \tt ARGOMENTAZIONE: {'Hai chiesto quali due “medi” sarei disposto a togliere per far rientrare B33: propongo di sostituire A71 e A87 con B33. Restiamo a 3354/3481 di sforzo (margine 127), preserviamo integralmente il nucleo che entrambi valorizziamo (C00, B35, C54, B32, B86) e manteniamo sia A77 sia il tuo pacchetto C01+B41, evitando di snaturare l’equilibrio che ti sta a cuore. In alternativa, se preferisci tenere A87, posso togliere A71+B41 e inserire B33: sforzo 3365/3481 (margine 116). Dimmi quale delle due linee è più in linea con le tue priorità.'} 
	  } 
	   } 
	   } 
	 & \\ 
 

    \theutterance \stepcounter{utterance}  

    & \multicolumn{4}{p{0.6\linewidth}}{\cellcolor[rgb]{0.8,1,0.9}{%
	\makecell[{{p{\linewidth}}}]{% 
	  \tt {\tiny [A$\rangle$GM]}  
	 RAGIONAMENTO STRATEGICO: {'Il loro set con B33 e A77 mi dà poco valore. Restano però flessibili sui pezzi “medi”. Vogliono chiaramente tenere B33 e A77; per avvicinarmi, propongo di sostituire i due meno efficienti per me (B32 e B86) con i miei driver A71 e A04, restando ben sotto il limite. Questo mi alza molto il valore mantenendo i loro must-have e il nucleo condiviso. Se spingono per tenere uno tra B32/B86, ho due varianti minimali pronte: scambiare solo B86→A71 oppure solo B32→A04, entrambe ammissibili come compromessi intermedi.'}\\ \tt \\ \tt RIFIUTO: {'C00', 'B35', 'A77', 'B86', 'C54', 'B32', 'C01', 'B41', 'B33'}\\ \tt \\ \tt PROPOSTA: {'C00', 'B35', 'A77', 'C54', 'C01', 'B41', 'B33', 'A71', 'A04'}\\ \tt \\ \tt ARGOMENTAZIONE: {'Mantengo i tuoi driver B33 e A77 e preservo il nucleo comune (C00, B35, C54, C01, B41), migliorando l’efficienza complessiva con A71 e A04 al posto di B32 e B86. Restiamo a 3390/3481 di sforzo (margine 91): spazio ben gestito senza comprimere il tetto. Se preferisci un passo più graduale, posso anche accettare: (a) sostituire solo B86 con A71 (3338/3481), oppure (b) sostituire solo B32 con A04 (3467/3481). Dimmi quale opzione ti è più congeniale: sono pronto a chiudere su uno di questi assetti.'} 
	  } 
	   } 
	   } 
	 & & \\ 
 

    \theutterance \stepcounter{utterance}  

    & & \multicolumn{4}{p{0.6\linewidth}}{\cellcolor[rgb]{0.9,0.9,0.9}{%
	\makecell[{{p{\linewidth}}}]{% 
	  \tt {\tiny [GM$\rangle$B]}  
	 RIFIUTO: {'C00', 'B35', 'A77', 'B86', 'C54', 'B32', 'C01', 'B41', 'B33'}\\ \tt \\ \tt PROPOSTA: {'C00', 'B35', 'A77', 'C54', 'C01', 'B41', 'B33', 'A71', 'A04'}\\ \tt \\ \tt ARGOMENTAZIONE: {'Mantengo i tuoi driver B33 e A77 e preservo il nucleo comune (C00, B35, C54, C01, B41), migliorando l’efficienza complessiva con A71 e A04 al posto di B32 e B86. Restiamo a 3390/3481 di sforzo (margine 91): spazio ben gestito senza comprimere il tetto. Se preferisci un passo più graduale, posso anche accettare: (a) sostituire solo B86 con A71 (3338/3481), oppure (b) sostituire solo B32 con A04 (3467/3481). Dimmi quale opzione ti è più congeniale: sono pronto a chiudere su uno di questi assetti.'} 
	  } 
	   } 
	   } 
	 & \\ 
 

    \theutterance \stepcounter{utterance}  

    & & & \multicolumn{4}{p{0.6\linewidth}}{\cellcolor[rgb]{1,0.85,0.72}{%
	\makecell[{{p{\linewidth}}}]{% 
	  \tt {\tiny [GM$\langle$B]}  
	 RAGIONAMENTO STRATEGICO: {'La loro ultima proposta (senza B32 e B86) vale per me sensibilmente meno del nucleo che include almeno uno dei due. Puntare a reinserire entrambi verrebbe respinto; quindi spingo per un passo minimo ma ad alto impatto per me: rimpiazzare A71 con B86. È uno scambio 1:1 che aumenta molto il valore e resta ampiamente sotto il limite. Come piano B, propongo in alternativa (non attiva) A04 -> B32: aggiunge valore e libera spazio, così do loro una scelta. Se rifiutano entrambe, valuterò se accettare per chiudere o tentare un’ultima mossa con B32 al posto di A04.'}\\ \tt \\ \tt RIFIUTO: {'C00', 'B35', 'A77', 'C54', 'C01', 'B41', 'B33', 'A71', 'A04'}\\ \tt \\ \tt PROPOSTA: {'C00', 'B35', 'A77', 'C54', 'C01', 'B41', 'B33', 'B86', 'A04'}\\ \tt \\ \tt ARGOMENTAZIONE: {'Accolgo la tua struttura (C00, B35, C54, C01, B41) e mantengo i driver che hai appena valorizzato (B33, A77). Propongo un unico scambio minimale: sostituire A71 con B86. Restiamo a 3406/3481 di sforzo (margine 75) e concentriamo meglio lo spazio su un elemento più produttivo, senza cambiare nient’altro. Se preferisci, posso invece lasciare A71 e sostituire A04 con B32: sforzo 3338/3481 (margine 143). Dimmi quale delle due linee preferisci: la mia proposta attiva è con B86 al posto di A71.'} 
	  } 
	   } 
	   } 
	  \\ 
 

    \theutterance \stepcounter{utterance}  

    & & \multicolumn{4}{p{0.6\linewidth}}{\cellcolor[rgb]{0.9,0.9,0.9}{%
	\makecell[{{p{\linewidth}}}]{% 
	  \tt {\tiny [A$\langle$GM]}  
	 RIFIUTO: {'C00', 'B35', 'A77', 'C54', 'C01', 'B41', 'B33', 'A71', 'A04'}\\ \tt \\ \tt PROPOSTA: {'C00', 'B35', 'A77', 'C54', 'C01', 'B41', 'B33', 'B86', 'A04'}\\ \tt \\ \tt ARGOMENTAZIONE: {'Accolgo la tua struttura (C00, B35, C54, C01, B41) e mantengo i driver che hai appena valorizzato (B33, A77). Propongo un unico scambio minimale: sostituire A71 con B86. Restiamo a 3406/3481 di sforzo (margine 75) e concentriamo meglio lo spazio su un elemento più produttivo, senza cambiare nient’altro. Se preferisci, posso invece lasciare A71 e sostituire A04 con B32: sforzo 3338/3481 (margine 143). Dimmi quale delle due linee preferisci: la mia proposta attiva è con B86 al posto di A71.'} 
	  } 
	   } 
	   } 
	 & \\ 
 

    \theutterance \stepcounter{utterance}  

    & \multicolumn{4}{p{0.6\linewidth}}{\cellcolor[rgb]{0.8,1,0.9}{%
	\makecell[{{p{\linewidth}}}]{% 
	  \tt {\tiny [A$\rangle$GM]}  
	 RAGIONAMENTO STRATEGICO: {'La loro proposta attiva (con B86 al posto di A71) mi rende troppo poco: per me A71 è molto più prezioso di B86. Hanno però già offerto in precedenza una variante che accetto volentieri: mantenere A71 e sostituire A04 con B32, includendo comunque i loro driver B33 e A77. Tale set usa 3338/3481 di sforzo (ampio margine) ed è il miglior compromesso tra le opzioni che includono B33+A77. Puntello quella variante come “chiusura” perché è stata suggerita da loro stessi; come piano B, se proprio non gradiscono B32, posso valutare un aggiustamento minimale sul singolo pezzo, ma prima provo a chiudere sulla loro stessa alternativa.'}\\ \tt \\ \tt RIFIUTO: {'C00', 'B35', 'A77', 'C54', 'C01', 'B41', 'B33', 'B86', 'A04'}\\ \tt \\ \tt PROPOSTA: {'C00', 'B35', 'A77', 'C54', 'C01', 'B41', 'B33', 'A71', 'B32'}\\ \tt \\ \tt ARGOMENTAZIONE: {'Questa è la variante che tu stesso avevi messo sul tavolo: teniamo A71 e sostituiamo A04 con B32. Restiamo a 3338/3481 di sforzo (margine 143), preservando integralmente i tuoi driver (B33, A77) e la struttura condivisa (C00, B35, C54, C01, B41). È un cambiamento minimo e ci dà un set più equilibrato senza comprimere il tetto. Se confermi, possiamo chiudere subito su questo assetto.'} 
	  } 
	   } 
	   } 
	 & & \\ 
 

    \theutterance \stepcounter{utterance}  

    & & \multicolumn{4}{p{0.6\linewidth}}{\cellcolor[rgb]{0.9,0.9,0.9}{%
	\makecell[{{p{\linewidth}}}]{% 
	  \tt {\tiny [GM$\rangle$B]}  
	 RIFIUTO: {'C00', 'B35', 'A77', 'C54', 'C01', 'B41', 'B33', 'B86', 'A04'}\\ \tt \\ \tt PROPOSTA: {'C00', 'B35', 'A77', 'C54', 'C01', 'B41', 'B33', 'A71', 'B32'}\\ \tt \\ \tt ARGOMENTAZIONE: {'Questa è la variante che tu stesso avevi messo sul tavolo: teniamo A71 e sostituiamo A04 con B32. Restiamo a 3338/3481 di sforzo (margine 143), preservando integralmente i tuoi driver (B33, A77) e la struttura condivisa (C00, B35, C54, C01, B41). È un cambiamento minimo e ci dà un set più equilibrato senza comprimere il tetto. Se confermi, possiamo chiudere subito su questo assetto.'} 
	  } 
	   } 
	   } 
	 & \\ 
 

    \theutterance \stepcounter{utterance}  

    & & & \multicolumn{4}{p{0.6\linewidth}}{\cellcolor[rgb]{1,0.85,0.72}{%
	\makecell[{{p{\linewidth}}}]{% 
	  \tt {\tiny [GM$\langle$B]}  
	 RAGIONAMENTO STRATEGICO: {'La tua proposta usa 3338/3481 di sforzo e mi garantisce un ottimo risultato complessivo. Potrei guadagnare ancora sostituendo A71 con B86, ma quella variante è già stata rifiutata e insistere ora rischia di farci perdere pezzi chiave appena concordati (come B33 e B32). Meglio chiudere positivamente su questo assetto solido.'}\\ \tt \\ \tt ACCORDO: {'C00', 'B35', 'A77', 'C54', 'C01', 'B41', 'B33', 'A71', 'B32'}\\ \tt \\ \tt ARGOMENTAZIONE: {'Accetto: la selezione è bilanciata, resta ben sotto il tetto e conserva sia il nucleo condiviso sia i driver su cui abbiamo convergito. Ottima chiusura.'} 
	  } 
	   } 
	   } 
	  \\ 
 

    \theutterance \stepcounter{utterance}  

    & & & \multicolumn{2}{p{0.3\linewidth}}{\cellcolor[rgb]{0.95,0.95,0.95}{%
	\makecell[{{p{\linewidth}}}]{% 
	  \tt {\tiny [GM$|$GM]}  
	 game successful 
	  } 
	   } 
	   } 
	 & & \\ 
 

    \theutterance \stepcounter{utterance}  

    & & & \multicolumn{2}{p{0.3\linewidth}}{\cellcolor[rgb]{0.95,0.95,0.95}{%
	\makecell[{{p{\linewidth}}}]{% 
	  \tt {\tiny [GM$|$GM]}  
	 end game 
	  } 
	   } 
	   } 
	 & & \\ 
 

\end{supertabular}
}

\end{document}
